\section{Related Works}
\label{sec:related}

\subsection{End-to-end Autonomous Driving}
End-to-end autonomous driving aims to directly generate motion planning from raw sensor inputs and has attracted increasing research attention~\cite{review1, review2, hagedorn2024integration, hu2025vision}. Most end-to-end autonomous driving systems follow a general framework that integrates perception, prediction, and planning in a cascaded~\cite{uniad, fusionad} or parallel manner~\cite{paradrive, transfuser} with a structured BEV feature~\cite{bevformer, li2024ego, diffusiondrive, gui2025trajdiff}. To reduce the reliance on dense BEV features, several works have proposed the use of sparse representations~\cite{vad, sparsead, sparsedrive, drivetransformer}. Leveraging driving priors, VADv2~\cite{vadv2} introduced a trajectory vocabulary as a prior for trajectory sampling. Building on this, Hydra-MDP~\cite{hydra} further distilled rule-based information. Additionally, DiffusionDrive~\cite{diffusiondrive} proposed a trajectory diffusion framework based on the trajectory vocabulary to accelerate the denoising process, and WoTE~\cite{li2025wote} utilized the trajectory anchor and future BEV state prediction to enhance the driving performance.

% TODO
\begin{figure*}[t]
    \centering
    \includegraphics[width=1.0\linewidth]{imgs/main_fig_v29.pdf}
    \caption{\textbf{Overall architecture of WorldDrive}. WorldDrive is a holistic framework unifying vision and motion representation to bridge scene generation and planning. The training process includes~Phase 1: WorldDrive for scene generation and Phase 2: WorldDrive for motion planning. The vision and motion representations are optimized through the scene generation task. In the planning stage, the planner utilizes the frozen vision and trajectory encoders and outputs top-$K$ multi-modal trajectories. A future-aware rewarder is further designed to select the optimal trajectory from the candidates.}
    \label{fig:main}
\end{figure*}

\subsection{Driving World Models}

Driving world models aim to predict the scene evolution from observations~\cite{review4, review5, review6, review7} and have shown great potential for generating long-horizon~\cite{magicdrivev2, vista, infinitydrive}, high-fidelity~\cite{drivedreamer, uniscene, drivingsphere}, and corner-case data~\cite{cosmos, terasim, sun2025terasim}. Although these methods are capable of generating driving scenarios, driving world models need to simulate reasonable scenarios based on motion conditions~\cite{jia2023adriver, hu2023gaia, genad, hassan2025gem}. Vista~\cite{vista} achieved strong dynamic modeling by utilizing a larger volume of driving scenario data and introducing a motion loss. Building on this, DriVerse~\cite{driverse} realized trajectory-specific video generation by encoding trajectories as textual prompts and motion priors, and incorporated a motion alignment module. ReSim~\cite{yang2025resim} enriched real-world human demonstrations with diverse non-expert data collected from a driving simulator. DrivingGPT~\cite{drivinggpt} and Epona~\cite{epona} introduced a discrete action representation at each timestep on top of an autoregressive video generation framework, enabling controllable trajectory generation.

\subsection{World Model for Planning}
Research on optimizing planning policy using the world model has begun to show potential due to the strong representation capability~\cite{review4, review5, review8}. Drive-WM~\cite{drivewm} was the first work to introduce the driving world model into end-to-end planning. It used predicted trajectories as conditions for future scenario prediction and combined perception modules to evaluate scenario safety. FSDrive~\cite{zeng2025fsdrive} proposed a visual chain-of-thought pipeline for future scene generation and planning. DrivingGPT~\cite{drivinggpt} and Epona~\cite{epona} used discrete and continuous autoregressive models to unify the generation of driving scenarios and driving trajectories, respectively. PWM~\cite{zhao2025pwm} and DriveVLA-W0~\cite{drivevlaw0} coupled trajectory planning with world simulation via action-free future state forecasting. Considering the high cost and latency of scene generation, latent world model methods, such as LAW~\cite{law}, World4Drive~\cite{world4drive}, and WorldRFT~\cite{yang2025worldrft}, demonstrated that effective scene representation can still be learned via self-supervision in latent space.